\scanpage{130}%
% archive sheet number: 129
% source lines 1--2; completes préliminaire from scan 129
minaire, - sous réserve tout au plus qu'il existe un tel $\unsure{\mathbb{E}|\mathbb{F}|E|F}$ sur une extension convenable de $k_x$ (chose peut-être toujours vrai\handwritten{e}).

% source lines 3--12
12. \uline{Exemples instructifs}. Soit $G$ un monoïde fini. Contrairement au yoga galoisien général, au lieu de prendre $A=k(G)$ (anneau des fonctions sur l'ensemble sous-jacent à $G$) prenons $U=k(G)$, avec l\deleted{a}'application diagonale déduite de $GxG\handwritten{\longrightarrow}G$ (la composition). (On a la bialgèbre ``\unsure{duale|daale}''\inserted{au sens Cartier} de celle qu'on considère d'habitude). Donc la catégorie $C$ définie par $G$ \deleted{est} peut être considérée comme la catégorie des familles $(\unsure{\mathbb{M}|\mathbb{X}}_g)_{g\ G}$ de vectoriels indexés par $G$, \unsure{lé|le}\deleted{application diagonale} produit\inserted{tens.} de deux tels $X$ et $Y$ étant $Z$ défini par $Z_g=\handwritten{\displaystyle\sum}_{g'g''=g}X_{g'}\otimes Y_{g''}$. L'unité est la ``famille de Dirac'' définie par $\mathbb{M}_g=0$ si $g\handwritten{\ne}\mathbb{e}$, $\mathbb{M}_e=k$.\inserted{On a une \unsure{Catégorie|catégorie} avec $\otimes,\handwritten{\phi},1$, et de plus un $\handwritten{\psi}$ si $G$ est commutatif, un $\vee$ si $G$ est un groupe.}

% source lines 13--24, with one displayed formula and its right marginal note
Un $L$ avec une application diagonale $L\handwritten{\longrightarrow}L\otimes_k L$ compatible avec $U\handwritten{\to}U\otimes_k U$ et donnant un isomorphisme par extensions des anneaux de base, peut s'interpréter comme une famille $(\deleted{M}\handwritten{L}_g)_{g\handwritten{\in}G}$, avec des isomorphismes donnés
\[
\deleted{M}\handwritten{L}_{gg'}\handwritten{\simeq}\deleted{M}\handwritten{L}_g\otimes\deleted{M}\handwritten{L}_{g'}\ .
\tag*{(*)}
\]
\margin{Ceci implique que le rang des $\deleted{M}\handwritten{L}_g$ est $0$ ou $1$}
C'est compatible avec $1$ si on a un isomorphisme\inserted{\handwritten{\unsure{(**)|(*+)}}} $\deleted{M}\handwritten{L}_e\handwritten{\simeq}k$, qui rendre les isomorphismes (*) triviaux quand \deleted{\unsure{ceux-ci}} $g$ ou $g'$ est $=1$. C'est associatif si et seulement si les isomorphismes (*) satisfont la \deleted{\unsure{conditio}} conditions d'associativité évidente. De même, $G$ étant commutatif, pour la commutativité. - Le foncteur $F'$ est toujours exact, car $M$ est plat sur $U$.

% source lines 25--30
a) Exemple d'un $\otimes$-foncteur exact, compatible avec $\handwritten{\phi,\psi},1$, \uline{mais non fidèle} (donc pas loc. \unsure{isoM.|isom.} au $\otimes$-foncteur donné). Ce n'est possible dans notre exemple que si $G$ n'est pas un groupe (\deleted{\unsure{sinon $N$|sinon $M$}} car $\deleted{M}\handwritten{L}_g$ est de rang $1$ pour $g$ inversible). Mais prenons $G=(e,p)$, avec $p^{\handwritten{2}}=p$, prenons $\deleted{M}\handwritten{L}_e=k$, $\deleted{M}\handwritten{L}_p=0$. On n'a plus le choix pour les isomorphismes $*$, qui satisfont toutes les compatibilités (avec $1_x$ par construction, $\handwritten{\phi,\psi}$).
% bottom margin, below a handwritten separating rule
\margin{\handwritten{\uline{Question} Un $\otimes$-foncteur compatible avec tout entre $\otimes$-catégories satisfaisant tous les axiomes (y compris rigidité) est-il néc. fidèle et exact ?!}}
